
\documentclass{article} % For LaTeX2e
\usepackage[submission]{colm2026_conference}

\usepackage{microtype}
\usepackage{hyperref}
\usepackage{url}
\usepackage{booktabs}
\usepackage{amsmath}
\usepackage{amssymb}
\usepackage{graphicx}

% NOTE: including geometry package
% The geometery package modifies some page properties when used. This can dramatically change the page margins, leading to severe template violation, and potential desk rejection. If the package is required, it can be used with the "pass" flag to skip the default page modifications, as in the following line:
% \usepackage[pass]{geometry}

\usepackage{lineno}

\definecolor{darkblue}{rgb}{0, 0, 0.5}
\hypersetup{colorlinks=true, citecolor=darkblue, linkcolor=darkblue, urlcolor=darkblue}

% Custom commands for notation
\newcommand{\zt}{z_t}
\newcommand{\ztnext}{z_{t+1}}
\newcommand{\rt}{r_t}

\title{DLE-Agent: A Dual-Level Evolution Framework for Stateful AI Research Agents}

% Authors must not appear in the submitted version. This should be be taken care of automatically as long as you are using the "submission" option for the colm2026_conference package. But it's on the authors to verify. Non-anonymous submissions will be rejected without review.

\author{Anonymous Authors}

% The \author macro works with any number of authors. There are two commands
% used to separate the names and addresses of multiple authors: \And and \AND.
%
% Using \And between authors leaves it to \LaTeX{} to determine where to break
% the lines. Using \AND forces a linebreak at that point. So, if \LaTeX{}
% puts 3 of 4 authors names on the first line, and the last on the second
% line, try using \AND instead of \And before the third author name.

\newcommand{\fix}{\marginpar{FIX}}
\newcommand{\new}{\marginpar{NEW}}

\begin{document}

\ifcolmsubmission
\linenumbers
\fi

\maketitle

\begin{abstract}
AI research agents have shown strong promise in automating machine learning
engineering, yet most existing systems remain cognitively stateless: they search
over candidate code solutions without maintaining or evolving an internal
understanding of the task. As a result, they often rely on shallow
trial-and-error, accumulate little task-specific knowledge across steps, and
struggle to explore efficiently on challenging multi-step problems. To address
this limitation, we propose \textbf{Dual-Level Evolution Agent (DLE-Agent)}, a
stateful research agent framework in which the agent not only searches for
better solutions but also builds and refines an explicit model of what it has
learned about the task. DLE-Agent operates through a dual-level mechanism. At
the outer level, the agent manages multiple research branches in a search tree,
adaptively deciding which branch to expand, continue, or abandon. At the inner
level, each branch maintains a structured cognitive state that encodes the
agent's current task understanding, working hypotheses, accumulated failure
patterns, and self-assessed confidence. This cognitive state is updated after
each execution step based on multi-dimensional feedback, including performance
trends, error types, and solution novelty. The updated state then shapes how
the agent generates and refines solutions in subsequent steps, forming a closed
loop where exploration at the branch level and learning at the cognitive level
continuously reinforce each other. We evaluate DLE-Agent on AIRS-Bench across
20 cross-domain tasks. Under identical step budgets and model configurations,
DLE-Agent with the cognitive layer achieves an average normalized score of
0.67, compared to 0.40 for the strongest stateless baseline, representing a
68\% relative improvement. This gain is consistent across three different
search strategies. Causal intervention experiments further confirm that the
cognitive state carries task-specific information and causally improves
downstream decisions, with the largest gains on complex multi-step tasks.
\end{abstract}
 


\section{Introduction}
\label{sec:intro}

Recent advances in large language models have enabled a new class of
systems---\emph{AI research agents}---that can read task specifications, write
and revise code, run experiments, and iteratively improve solutions with
limited human intervention~\citep{huang2024mlagentbench, chan2024mlebench,
airadojo2025, airsbench2025}. These agents are particularly promising for
machine learning research tasks, where solving a problem typically requires
multiple rounds of hypothesis formation, implementation, diagnosis, and
revision rather than one-shot code generation. In this setting, the central
challenge is not only to produce candidate solutions, but also to
progressively build and update task understanding over long-horizon
interaction.

Despite strong progress, many existing research agents still lack an explicit,
structured, and continuously updated representation of what has been learned so
far about the task~\citep{shinn2023reflexion, zhao2024expel,
wang2023voyager}. Prior information is often retained only implicitly---through
dialogue history, execution logs, local reflections, or compressed
summaries. Such mechanisms help preserve context, but they do not amount to a
persistent internal state that can be directly maintained, revised, and reused
across decision steps. As a result, current agents can often \emph{access}
past interactions without reliably \emph{accumulating understanding} from them.

This limitation leads to three recurring failure modes. First, agents tend to
\textbf{repeat similar mistakes}---such as revisiting known implementation
bugs, environment mismatches, or invalid directions---thereby wasting limited
search budget. Second, they often \textbf{fail to consolidate} task-specific
knowledge acquired in earlier attempts, including observations about promising
model families, stable training regimes, or unproductive regions of the search
space. Third, their exploration is frequently \textbf{myopic and unstable}:
without an explicit representation of what has been tried and what remains
plausible, agents tend to oscillate between directions instead of refining a
coherent research trajectory.

We address this problem with \textbf{D}ual-\textbf{L}evel
\textbf{E}volution \textbf{Agent} (\textbf{DLE-Agent}), a framework that treats AI research as a
closed loop between search and cognitive-state evolution. Instead of using past
experience solely as passive context, DLE-Agent maintains, for each research
trajectory, a structured \emph{cognitive state} $z_t$ that explicitly
represents the agent's current task understanding, active hypotheses, learned
failure patterns, directional preferences, and confidence about the current
line of investigation. After each execution step, this state is updated
through a \emph{reflection} process that incorporates feedback such as
performance changes, error types, and trial novelty. The updated cognitive
state is then injected into subsequent proposal, improvement, and debugging
decisions, allowing accumulated understanding to directly shape future actions.

From a systems perspective, DLE-Agent operates on two coupled levels: an
\emph{outer level} that determines how search proceeds over research
trajectories---deciding which branch to expand, continue, or
abandon---and an \emph{inner level} that continuously refines the agent's
internal task representation. The key question, therefore, is not simply
whether additional memory helps, but whether an explicit and evolving
cognitive state can systematically improve downstream research behavior.
DLE-Agent can be instantiated with different outer-level search
structures---from chain-based sequential iteration to multi-draft parallel
exploration to tree-based search---providing a controlled experimental
platform for studying \emph{when and why} explicit cognitive state is
effective.

To study this question, we evaluate DLE-Agent on
AIRS-Bench~\citep{airsbench2025}, a benchmark of research-oriented machine
learning tasks spanning multiple domains. Our goal is not only to compare
final task performance, but also to test whether the cognitive state itself
carries task-relevant information and causally influences later decisions. To
this end, we design controlled state interventions---including
\emph{ablation} (blanking the state), \emph{freezing} (halting evolution at an
early step), and \emph{cross-task replacement} (substituting a state from an
unrelated task)---to distinguish the effect of cognitive-state evolution from
the superficial benefit of adding more textual context. Our results suggest
that explicit cognitive-state evolution is particularly helpful for tasks that
require iterative diagnosis, multi-step reasoning, and knowledge accumulation,
and less impactful when progress is dominated by external computational
constraints.

More broadly, this work frames a research question for agent design: when an
agent can explicitly maintain and update its own task understanding, under
what conditions does this improve its behavior, and why? We hope this
perspective can help move research agents from repeated trial-and-error toward
more stateful and cumulative forms of exploration.

\paragraph{Contributions.} Our main contributions are:
\begin{itemize}
  \item We propose \textbf{DLE-Agent}, a framework that couples search over
  research trajectories with explicit cognitive-state evolution, making task
  understanding a first-class object in the agent's decision loop.
  \item We introduce a \textbf{structured cognitive-state representation} and
  a \textbf{reflection-based update mechanism} for accumulating task
  understanding, working hypotheses, failure patterns, and directional
  preferences across steps.
  \item We develop a \textbf{controlled experimental protocol} on AIRS-Bench
  to study not only \emph{whether} cognitive state helps, but also
  \emph{when} and \emph{why} it helps, through systematic state
  interventions.
  \item We provide \textbf{initial evidence} that explicit cognitive-state
  evolution is most beneficial for tasks requiring sustained reasoning,
  diagnosis, and knowledge accumulation, offering empirical grounding for
  the question of when thinking helps.
\end{itemize}

% ====================================================================
% Section 2: Related Work
% ====================================================================
\section{Background and Related Work}
\label{sec:related}

\subsection{AI research agents and internal state}
\label{sec:rw-agents}

The dominant paradigm for AI research agents treats ML engineering as search
over candidate code solutions. AIDE~\citep{jiang2025aide} pioneered this view
by searching directly in code space via Draft, Debug, and Improve operators.
AIRA-dojo~\citep{airadojo2025} unified multiple search policies (greedy,
MCTS, evolutionary) within a single framework and showed that \emph{operator
design is a stronger bottleneck than search policy}. In both systems, however,
the search state consists solely of external artifacts---candidate programs
and their scores---while the understanding the LLM develops during search
remains implicit in the transient context window.

Several recent systems introduce persistent internal state to address this
limitation. ML-Master~2.0~\citep{zhu2026mlmaster} uses hierarchical cognitive
caching to distill execution traces into progressively stable knowledge, but
its memory is fundamentally passive storage that records \emph{what happened}
without representing what the agent should believe or do next.
R\&D-Agent~\citep{yang2025rdagent} separates research from development and
evolves hypotheses across iterations; however, its hypotheses are high-level
natural-language statements rather than structured, multi-field state objects
updated after every step. LoongFlow~\citep{wan2025loongflow} combines
evolutionary search with a Plan-Execute-Summarize paradigm, but its
experience memory operates at the population level rather than per-trajectory.
In the broader agent literature, Reflexion~\citep{shinn2023reflexion}
cumulatively appends verbal reflections to the prompt, yet these are
unstructured text subject to context-window truncation.
Self-Refine~\citep{madaan2023selfrefine} refines within a single step without
cross-step accumulation; LATS~\citep{zhou2024lats} uses reflections as
auxiliary value estimates during tree search;
Voyager~\citep{wang2023voyager} stores successful behaviors as reusable code
snippets---passive storage rather than a cognitive representation of task
understanding.

A common thread across these systems is that internal state, when present, is
evaluated only via component ablation: removing the module and observing
performance drop. This conflates the value of having \emph{any} state
structure with the value of its task-specific \emph{content}. DLE-Agent
addresses this gap through state-intervention experiments (e.g., replacing the
evolved state with one from an unrelated task), directly testing whether the
content---not merely the mechanism---drives the observed benefit.

\subsection{Positioning}
\label{sec:rw-position}

Table~\ref{tab:comparison} compares DLE-Agent with representative systems
along four design dimensions.

\begin{table}[t]
\centering
\caption{Design-dimension comparison. All systems search in code space;
columns describe \emph{additional} internal-state mechanisms.}
\label{tab:comparison}
\footnotesize
\renewcommand{\arraystretch}{1.15}
\begin{tabular}{@{}l@{\hskip 6pt}l@{\hskip 6pt}l@{\hskip 6pt}l@{\hskip 6pt}l@{}}
\toprule
\textbf{System}
  & \textbf{Internal state}
  & \textbf{Persist.}
  & \textbf{Update}
  & \textbf{State role} \\
\midrule
AIDE / AIRA-dojo
  & Artifact history
  & Prompt win.
  & Concatenation
  & Search context \\
ML-Master~2.0
  & Hier.\ cache
  & Cross-step
  & Rule migration
  & Passive store \\
R\&D-Agent
  & Evolving hyp.
  & Cross-iter.
  & Phase switch
  & Guide research \\
LoongFlow
  & Evol.\ memory
  & Population
  & Evol.\ selection
  & Shared pool \\
Reflexion
  & Verbal refl.
  & Cross-attempt
  & LLM append
  & Prompt augment. \\
\midrule
\textbf{DLE-Agent}
  & \textbf{Structured $z_t$}
  & \textbf{Cross-step}
  & \textbf{LLM reflect}
  & \textbf{Active shaping} \\
\bottomrule
\end{tabular}
\end{table}

DLE-Agent's contribution is not in ``adding a cognitive layer''---a direction
explored independently by multiple teams---but in two aspects: (i)~providing
\emph{causal-level evidence} via state-intervention experiments that verify
the information value of cognitive-state \emph{content} rather than structure,
and (ii)~systematically analyzing the \emph{boundary conditions} of cognitive
evolution---on which tasks it helps, on which it does not, and when the
overhead is justified.


% ====================================================================
% Section 3: Method
% ====================================================================
\section{Method}
\label{sec:method}

\subsection{Problem formulation}
\label{sec:formulation}

Following AIRA-dojo~\citep{airadojo2025}, we formalize an AI research agent as
a tuple operating within an environment. AIRA-dojo defines the agent as a
3-tuple $(\pi, \mathcal{O}, \mathcal{E})$: a search policy~$\pi$, an operator
set~$\mathcal{O}$, and an execution environment~$\mathcal{E}$. At each
step~$t$, the policy selects an operator $o \in \mathcal{O}$, which generates a
code solution (a \emph{node}); the environment executes the code and returns a
scalar metric. In this framework, the search object is the external code
artifact, and the agent's ``knowledge'' lives implicitly in the prompt context,
neither explicitly maintained nor reused across steps.

DLE-Agent extends this formulation to a 5-tuple
$(\pi, \mathcal{O}, \mathcal{E}, z_0, U)$ by introducing two new components:
\begin{itemize}
  \item \textbf{Initial cognitive state}~$z_0$: a structured representation
  initialized at the start of each task, encoding the agent's preliminary
  understanding (Section~\ref{sec:cognitive_state}). Subsequent states
  $z_1, z_2, \ldots$ are produced recursively by~$U$.
  \item \textbf{State update function}~$U$: a reflection operator that updates
  the cognitive state based on multi-dimensional execution feedback,
  $\ztnext = U(\zt, \rt)$ (Section~\ref{sec:inner}).
\end{itemize}

The per-step cycle becomes:
\begin{equation}
  o_t = \pi(\zt, \mathcal{J}_t) \;\longrightarrow\;
  a_t = o_t(\zt) \;\longrightarrow\;
  \rt = \mathcal{E}(a_t) \;\longrightarrow\;
  \ztnext = U(\zt, \rt)
  \label{eq:cycle}
\end{equation}
where $\mathcal{J}_t$ is the exploration journal (all nodes generated up to
step~$t$), $o_t \in \mathcal{O}$ is the selected operator, $a_t$ is the code
solution generated with~$\zt$ injected into the operator's prompt, and~$\rt$ is
the multi-dimensional feedback signal. Each of the five components appears
exactly once in the cycle.

When the cognitive-state module is disabled---that is, $z_0$ is not
initialized and $U$ is not invoked---the system reduces to standard AIRA-dojo.
This degeneracy provides a clean comparison: the two configurations share
identical~$\pi$, $\mathcal{O}$, and~$\mathcal{E}$; the only difference is the
presence or absence of the cognitive-state module.


\subsection{Cognitive state representation}
\label{sec:cognitive_state}

The cognitive state~$\zt$ is a structured object with eight semantic fields:
\begin{equation}
  \zt = \big(\,
  \underbrace{\tau}_{\text{task und.}},\;
  \underbrace{H}_{\text{hypotheses}},\;
  \underbrace{P}_{\text{patterns}},\;
  \underbrace{S}_{\text{attempts}},\;
  \underbrace{D^{+}}_{\text{preferred}},\;
  \underbrace{D^{-}}_{\text{avoided}},\;
  \underbrace{c}_{\text{confidence}},\;
  \underbrace{\varepsilon}_{\text{env ctx}}
  \,\big)
  \label{eq:state}
\end{equation}

\begin{itemize}
  \item \textbf{Task understanding}~$\tau$ (string): a natural-language summary
  of the agent's current interpretation of the task, refined as the agent
  discovers data properties, evaluation criteria, and domain constraints.
  \item \textbf{Hypotheses}~$H$ (string list): active working hypotheses about
  which approaches may succeed (e.g., ``GNN with edge features outperforms MLP
  on this molecular graph structure'').
  \item \textbf{Learned patterns}~$P$ (string list): failure and success
  patterns extracted from execution feedback (e.g.,
  ``\texttt{ReduceLROnPlateau(verbose)} removed in PyTorch~$\geq$\,2.4'').
  \item \textbf{Attempt summaries}~$S$ (record list): a compressed log of each
  attempt, storing the step number, approach description, metric value,
  bug status, and key insight. $S$~grows deterministically (one record appended
  per step) and is never rewritten by the LLM.
  \item \textbf{Preferred / avoided directions}~$D^{+}, D^{-}$ (string lists):
  directions the agent considers promising or has ruled out, forming explicit
  search constraints.
  \item \textbf{Confidence}~$c \in [0, 1]$: a self-assessed confidence in the
  current best approach; this value directly influences exploration--exploitation
  decisions in the MC-ESES search policy (Section~\ref{sec:outer}).
  \item \textbf{Environment context}~$\varepsilon$ (string): a structured
  summary of the execution environment---package versions, GPU capabilities,
  known API compatibility issues---populated by an initial reconnaissance script
  and carried forward unchanged across all subsequent steps.
\end{itemize}

\paragraph{Distinction from passive memory.}
The core fields of~$\zt$ ($\tau, H, P, D^{+}, D^{-}$) do not attempt to store
raw interaction history; that role is served by the attempt summaries~$S$
and the exploration journal~$\mathcal{J}_t$. Instead, these fields encode
\emph{inferential products}: ``the key challenge in this task is class
imbalance'' is not a datum but a judgment formed by integrating feedback across
multiple execution steps. This distinction parallels the difference between
remembering a phone number (passive storage) and understanding a proof
(active cognition).

\paragraph{Initialization.}
$z_0$ is produced by a single LLM call that takes the task description as
input. At this point, $\tau$~reflects an initial reading of the task,
$H$~contains prior-based guesses, $P = D^{+} = D^{-} = \varnothing$,
$S = \varnothing$, and $c$~is set to a low value indicating initial
uncertainty. The environment context~$\varepsilon$ is optionally populated
by a reconnaissance script that probes the execution environment for Python
package versions, GPU hardware, known API deprecations, and the data directory
structure.

\paragraph{Prompt injection.}
At each step, $\zt$ is rendered into a structured text block via a
\texttt{to\_prompt\_str()} method and injected in its entirety into the
current operator's prompt. All operators (draft, improve, debug) receive the
same complete rendering, which includes environment context, task
understanding, active hypotheses, learned patterns, preferred and avoided
directions, recent attempt summaries, and the confidence score.

\paragraph{Intrinsic quality.}
To measure the informational richness of a cognitive state independently of
task performance, we define a heuristic intrinsic quality score:
\begin{equation}
  \text{IQ}(\zt) = \frac{1}{|\mathcal{S}|}
  \sum\nolimits_{s \in \mathcal{S}} s
  \label{eq:iq}
\end{equation}
where $\mathcal{S}$ consists of the following sub-scores:
(i)~hypothesis richness, $\min(|H|/5,\, 1)$;
(ii)~pattern accumulation, $\min(|P|/5,\, 1)$;
(iii)~direction clarity,
$\tfrac{1}{2}[\min(|D^{+}|/3,\, 1) + \min(|D^{-}|/3,\, 1)]$;
and (iv)~confidence calibration, $\max(0,\, 1 - |c - \hat{c}|)$,
where $\hat{c}$ is the empirical non-buggy rate across all entries in~$S$.
The calibration term is included only when $S \neq \varnothing$, so
$|\mathcal{S}| \in \{3, 4\}$.
IQ is not a proxy for task score---a state that has accurately identified every
dead end may have high IQ but low task performance.
It is used as a secondary signal in MC-ESES node selection
(Section~\ref{sec:outer}).


\subsection{Execution-aware system scaffold}
\label{sec:scaffold}

Before cognitive-state evolution can be effective, the underlying agent must
be able to produce non-buggy, environment-compatible code. We introduce two
categories of execution-aware improvements that are shared by all experimental
conditions---including the ablated baseline---and therefore orthogonal to the
cognitive-state mechanism evaluated in Section~\ref{sec:causal}.

\paragraph{Environment-aware context construction.}
At the start of each run, a lightweight reconnaissance script is executed
inside the task environment. It probes the installed Python version, key
package versions (PyTorch, Transformers, scikit-learn, LightGBM, etc.), GPU
count and memory, and known API deprecations. The output is formatted as a
structured summary and injected into the environment context
field~$\varepsilon$ of the initial cognitive state~$z_0$
(Section~\ref{sec:cognitive_state}). Critically, $\varepsilon$ is preserved
even under ablation (ABL), so both NAT and ABL benefit equally from this
information. In addition, the package list reported to operators includes full
version specifiers (e.g., \texttt{torch==2.10.0} rather than just
\texttt{torch}), enabling the LLM to reason about API compatibility.

\paragraph{Failure-preserving execution interface.}
Two changes improve the quality of execution feedback returned to the agent.
First, operator prompts are augmented with environment-specific constraints
(e.g., that \texttt{ReduceLROnPlateau} no longer accepts \texttt{verbose} in
PyTorch~$\geq$\,2.4, that the \texttt{evaluate} library is unavailable, that
\texttt{fp16} should be preferred over \texttt{bf16} on the current hardware).
The full list of constraints is provided in Appendix~\ref{sec:app_prompts}.
Second, terminal output truncation is redesigned to preserve error
tracebacks: instead of retaining only the first and last $k$~characters
(which often discards errors appearing mid-output), the system first extracts
all \texttt{Traceback} and \texttt{Error} blocks, then fills the remaining
budget with head and tail context. This ensures the agent always sees the
full error when debugging.

These improvements address the most common failure modes of the original
AIRA-dojo scaffold. Their aggregate contribution is quantified in
Section~\ref{sec:decompose}.


\subsection{Dual-level evolution}
\label{sec:dual_level}

DLE-Agent's core design is \emph{dual-level evolution}: the agent
simultaneously evolves in cognitive space (inner level) and in the search tree
(outer level), with the two levels continuously informing each other. The inner
level defines how the cognitive state is updated from execution
feedback---this is the central contribution of the paper. The outer level
defines the search topology---reusing existing strategies as controlled
experimental variables.


\subsubsection{Inner level: cognitive state evolution}
\label{sec:inner}

After each code execution, the cognitive state is updated via a reflection
operator:
\begin{equation}
  \ztnext = U(\zt, \rt)
  \label{eq:reflect}
\end{equation}

\paragraph{Multi-dimensional feedback~$\rt$.}
Rather than conditioning only on a scalar metric and a binary success flag, we
construct a structured feedback signal with eight dimensions:
\begin{equation}
  \rt = \big(\,\text{metric},\;
  \text{is\_bug},\;
  \text{err\_cat},\;
  \text{err\_pat},\;
  \text{trend},\;
  \Delta,\;
  \text{nov},\;
  \Delta F\,\big)
  \label{eq:feedback}
\end{equation}
where $\text{err\_cat} \in
\{\text{environment, logic, data, resource, none}\}$ is determined by
rule-based pattern matching on terminal output (covering common patterns such
as CUDA OOM, missing modules, and deprecated APIs);
$\text{trend} \in \{\text{improving, stagnating, degrading, first}\}$ compares
the current metric to the historical best;
$\Delta = \text{metric}_t - \text{metric}_{\text{best}}$ is the signed gap;
$\text{nov} \in [0, 1]$ is the line-level Jaccard distance between the current
code and recent attempts;
and $\Delta F \in [0, 1]$ is a heuristic proxy for uncertainty reduction,
computed as a linear combination of novelty, trend, and bug status.
All dimensions are computed directly from execution results without additional
LLM calls.

\paragraph{Reflection operator.}
The update function~$U$ is implemented as a single LLM call. The LLM receives:
(i)~the current cognitive state~$\zt$ rendered as structured text;
(ii)~the executed code;
(iii)~the terminal output, with~$\rt$ appended as a formatted text block; and
(iv)~an analysis summary and the best metric so far.
The LLM is prompted to return a JSON object containing updated values for the
six mutable fields ($\tau, H, P, D^{+}, D^{-}, c$) together with a key
insight from the current attempt.

\paragraph{Update mechanism.}
The LLM's output for each field \emph{replaces} the corresponding field
in~$\zt$ to form~$\ztnext$. This is a full replacement, not an incremental
append: the LLM may freely add or remove hypotheses, reorganize patterns, and
shift directions. Two fields are carried forward deterministically and are
never modified by the LLM:
\begin{itemize}
  \item $\varepsilon$ (environment context): inherited unchanged from~$z_0$.
  \item $S$ (attempt summaries): a new record is appended after each step,
  containing the step number, approach summary, metric, bug status, and key
  insight. When $|S|$ exceeds a cap (default~50), the oldest entries are
  trimmed.
\end{itemize}

\paragraph{Fault tolerance.}
The LLM output is first parsed as JSON. If parsing fails or the result does not
contain the required \texttt{task\_understanding} field, the system falls back
to $\ztnext = \zt$ (no update), but still appends the current attempt to~$S$.
For successfully parsed responses, any omitted optional field inherits its
value from~$\zt$.

\paragraph{Reflection frequency.}
In the current implementation, reflection is invoked after every code
execution. In the Greedy topology this is controlled by a configuration flag
(\texttt{reflect\_after\_every\_step}); in MC-ESES, reflection runs after each
expansion child.


\subsubsection{Outer level: search topologies}
\label{sec:outer}

The inner-level cognitive state evolution composes orthogonally with different
outer-level search topologies. We use three topologies of increasing
complexity:

\paragraph{Chain ($K{=}1$).}
A single linear sequence: each step generates one solution, evaluates it,
reflects, and continues. There is no branching and no parallel drafts.
Chain is the simplest topology, isolating the contribution of cognitive-state
evolution: any performance gain must come from the interplay between the
operator's own refinement capability and the accumulated cognitive state, since
the search topology provides no diversity.

\paragraph{Greedy ($K{=}5$).}
Five independent drafts are generated in the first round, and the
highest-scoring draft is selected for iterative improvement. A rule-based
policy selects \textsc{debug} when a recent node is buggy, and
\textsc{improve} otherwise. The cognitive state influences code generation
through prompt injection but does not affect operator selection. This
topology matches AIDE~\citep{jiang2025aide} and the Greedy strategy in
AIRA-dojo~\citep{airadojo2025}.

\paragraph{MC-ESES ($K{\geq}1$, tree search).}
Monte Carlo Endogenous State Evolution Search---the full dual-level
optimization. MC-ESES operates in two phases.

\emph{Phase~1 (linear warm-up).} The agent accumulates an initial cognitive
state in a chain-like manner, but with a key difference: the search
policy~$\pi$ is directly driven by~$\zt$.
Specifically, if $c < 0.3$ or $D^{+}$ contains untried directions, the policy
selects \textsc{draft} (exploration); if all recent attempts are buggy and
debuggable nodes exist, it selects \textsc{debug}; if $c \geq 0.3$ and a
non-buggy node exists, it selects \textsc{improve} (exploitation).
This means that the cognitive state shapes not only the \emph{content} of
generated code (via prompt injection) but also the agent's \emph{macro-level
behavior} (via confidence-driven action selection).

\emph{Phase~2 (tree search).} The agent maintains a search tree in which
\textbf{each node carries an independent cognitive state}~$\zt^{(i)}$. Child
nodes are created by deep-copying the parent's cognitive state and evolving it
independently through execution and reflection. MC-ESES follows the standard
MCTS cycle (select $\to$ expand $\to$ evaluate $\to$ backpropagate) with a
modified selection rule.

Node selection uses a UCT formula that blends the backpropagated value with
intrinsic quality via a convex combination:
\begin{equation}
  \text{UCT}(i) = \big[(1 - w_{\text{iq}}) \cdot \bar{Q}_i
  + w_{\text{iq}} \cdot \text{IQ}(\zt^{(i)})\big]
  + c_{\text{uct}} \sqrt{\frac{\ln N}{n_i}}
  \label{eq:uct}
\end{equation}
where $\bar{Q}_i$ is the min-max normalized backpropagated value,
$N$ is the parent visit count, $n_i$ is the child visit count,
$c_{\text{uct}}$ controls the exploration--exploitation balance, and
$w_{\text{iq}}$ (default~0.2) controls the intrinsic quality weight.
The convex combination ensures that, for nodes with similar~$\bar{Q}$,
those with richer cognitive states are preferred.

Backpropagation uses a composite value:
\begin{equation}
  v = \alpha \cdot m
  + (\beta \cdot v_{\text{rate}} + \gamma \cdot I)\, |m|
  \label{eq:backprop}
\end{equation}
where $m$ is the best metric obtained in the expansion,
$v_{\text{rate}} = n_{\text{valid}} / n_{\text{total}}$ is the fraction of
non-buggy children,
$I \in \{0, 1\}$ indicates whether the expansion improved the global best,
and $|m|$ scales the auxiliary terms to align with the metric's magnitude.
The weights $\alpha, \beta, \gamma$ are read from configuration (defaults in
Section~\ref{sec:setup}).

\paragraph{Cross-branch knowledge transfer.}
When a leaf node is selected for expansion, the system harvests learned
patterns~$P$ and avoided directions~$D^{-}$ from the top-$k$ leaves (ranked by
$Q$-value, excluding the current node's ancestor chain). These insights are
rendered as an additional ``Cross-Branch Insights'' section in the operator
prompt, allowing knowledge to flow across branches without modifying any
node's cognitive state object or merging code.


\subsection{Causal intervention framework}
\label{sec:intervention}

To establish that the cognitive state \emph{causally} improves agent
behavior---rather than merely providing generic prompt augmentation---we
design a suite of controlled interventions. All conditions share the same
search topology, operator set, LLM, step budget, and random seed; the only
variable is the content of~$\zt$ injected into operator prompts.

Interventions are applied after each reflection step and before the next
execution---i.e., as post-processing on the output of $U(\zt, \rt)$:

\begin{table}[t]
\caption{Causal intervention conditions. All conditions share the same
search topology, operator set, LLM, step budget, and random seed;
only the content of~$\zt$ differs.}
\label{tab:interventions}
\centering
\footnotesize
\begin{tabular}{@{}lp{4.8cm}p{4.5cm}@{}}
\toprule
\textbf{Cond.} & \textbf{Treatment of $\zt$} & \textbf{Tests} \\
\midrule
NAT & No intervention; $\zt$ evolves normally.
    & Control condition. \\
ABL & Reset to blank state after each reflection
      (preserving~$\varepsilon$ and search-neutral fields).
    & Removes accumulated knowledge; retains prompt structure. \\
SCR & Replaced with terminal state from a completed run on an
      \emph{unrelated} task (preserving~$\varepsilon$).
    & Distinguishes task-specific content from generic enrichment. \\
FRZ & Evolves normally until step~$N_f$, then frozen.
    & Marginal value of continued evolution after initial learning. \\
\bottomrule
\end{tabular}
\end{table}

\paragraph{Implementation note.}
The precise behavior of ABL varies across topologies. In MC-ESES, where the
search policy~$\pi$ reads~$c$ and~$S$ from~$\zt$ (Section~\ref{sec:outer}),
ABL preserves the attempt summaries~$S$ and sets~$c = 0.5$ (a neutral value
that permits both draft and improve actions) to ensure that the
\emph{action distribution} matches NAT and only the \emph{prompt content}
differs. In Greedy, where~$\pi$ does not depend on~$\zt$, ABL preserves
only~$\varepsilon$. Both variants clear $\tau, H, P, D^{+}, D^{-}$,
ensuring a clean ablation of accumulated cognitive content.

\paragraph{Predicted ranking.}
If $\zt$ is merely a prompt filler, ABL and SCR should perform comparably to
NAT. If $\zt$ carries task-specific causal information, we expect
$\text{NAT} > \text{FRZ} > \text{ABL} \approx \text{SCR}$,
with FRZ approaching NAT as $N_f$ increases.


% ====================================================================
% Section 4: Experiments
% ====================================================================
\section{Experiments}
\label{sec:experiments}

We design three complementary analyses. First, we \emph{decompose the
performance gain} (\S\ref{sec:decompose}) into system-level scaffold
improvements and cognitive-state evolution, establishing what each component
contributes. Second, we perform a \emph{causal intervention analysis}
(\S\ref{sec:causal}) to test whether the cognitive state carries task-specific
information that causally improves downstream decisions. Third, we analyze
\emph{when cognitive state helps} (\S\ref{sec:boundary}), identifying
task-level boundary conditions. All analyses share the setup described below.


\subsection{Setup}
\label{sec:setup}

\paragraph{Benchmark and scoring.}
We evaluate on AIRS-Bench~\citep{airsbench2025}, which comprises 20~tasks
across five domains: natural language processing, code generation, mathematics,
time-series forecasting, and molecular property prediction. Each task provides
training data, a held-out test set, and an automatic evaluation metric. We
report the \emph{normalized score} (NS) defined by AIRS-Bench, which applies a
March-of-9s transform $\varphi(s) = -\log_{10}|s - s^{\text{opt}}|$ to map raw
scores onto a comparable scale where $\text{NS} = 0$ corresponds to the
estimated worst agent and $\text{NS} = 1$ corresponds to published SOTA.
Scores above~1 indicate that the agent surpasses SOTA. Per-task NS values are
averaged across seeds and then across tasks to yield the aggregate Avg~NS.

\paragraph{Task selection.}
We select 10~representative tasks spanning four difficulty tiers:
Easy (SICK Accuracy, Solar Weekly~MAE),
Medium (Yelp Accuracy, Winogrande Accuracy, SQuAD Exact Match),
Hard (Qm9~MAE, Zinc~MAE, Code Retrieval~MRR),
and Expert (SuperGLUE~WSC Accuracy, SVAMP Accuracy).
Both experiments use the same 10~tasks.

\paragraph{Model and infrastructure.}
All experiments use Qwen-3.5-397B-A17B as the backbone LLM, accessed via the
DashScope API~(Alibaba Cloud). Agent-generated code is executed as isolated
Python subprocesses on a server with 4$\times$NVIDIA RTX~5090 GPUs
(32\,GB each). The execution timeout is 4~hours per step. We use 3~random seeds
per condition. The step budget is 50~steps per seed for all topologies;
MC-ESES's tree-search branching operates within this same budget.

\paragraph{Hyperparameters.}
Table~\ref{tab:hyperparams} lists the key hyperparameters used across all
experiments.

\begin{table}[h]
\centering
\caption{Hyperparameters. Values are fixed across all tasks and seeds.}
\label{tab:hyperparams}
\footnotesize
\begin{tabular}{@{}llr@{}}
\toprule
\textbf{Component} & \textbf{Parameter} & \textbf{Value} \\
\midrule
Search (all)  & Step budget per seed      & 50 \\
              & Max debug depth            & 3 \\
Greedy        & Initial drafts $K$         & 5 \\
MC-ESES       & UCT exploration $c_{\text{uct}}$ & 1.41 \\
              & Intrinsic quality weight $w_{\text{iq}}$ & 0.2 \\
              & Cross-branch top-$k$       & 3 \\
Backprop      & Metric weight $\alpha$     & 0.6 \\
              & Validity weight $\beta$    & 0.2 \\
              & Improvement weight $\gamma$& 0.2 \\
Cognitive     & Attempt summary cap $|S|_{\max}$ & 50 \\
state         & Confidence draft threshold & 0.3 \\
Intervention  & FRZ freeze step $N_f$      & 5 \\
\bottomrule
\end{tabular}
\end{table}

\paragraph{Experimental conditions.}
\begin{itemize}
  \item \textbf{Cross-strategy comparison}: three topologies
  (Chain, Greedy, MC-ESES) $\times$ two conditions ($+\zt$: cognitive state
  enabled; $-\zt$: cognitive state module disabled).
  $10 \times 3 \times 6 = 180$~runs.
  \item \textbf{Causal intervention}: Greedy topology with four conditions
  (NAT, ABL, SCR, FRZ) as defined in Table~\ref{tab:interventions}.
  $10 \times 3 \times 4 = 120$~runs. The 30~Greedy~$+\zt$ runs are shared
  with the cross-strategy comparison (Greedy~$+\zt$ = NAT).
\end{itemize}

\paragraph{Baselines.}
We compare three conditions that form a layered attribution chain:
\begin{enumerate}
  \item \textbf{AIRA-dojo}~\citep{airadojo2025}: the unmodified Greedy
  scaffold from Meta's AIRA framework, using the original AIDE operator
  prompts without environment reconnaissance or cognitive state. We run
  AIRA-dojo with the same LLM (Qwen-3.5-397B), 50~steps, and 3~seeds
  for a matched comparison.
  \item \textbf{ABL} (DLE-Agent $-\zt$): DLE-Agent with the cognitive-state
  module set to \emph{ablated} mode---the reflection operator is invoked
  but its output is discarded, producing a blank state at every step.
  This preserves all system-level improvements
  (Section~\ref{sec:scaffold}) while removing cognitive-state
  accumulation.
  \item \textbf{NAT} (DLE-Agent $+\zt$): DLE-Agent with full
  cognitive-state evolution enabled.
\end{enumerate}
The chain Baseline~$\to$~ABL isolates the contribution of the execution-aware
scaffold; ABL~$\to$~NAT isolates the contribution of cognitive-state evolution.


\subsection{Decomposing system-level and state-evolution gains}
\label{sec:decompose}

Before analyzing cognitive-state evolution, we first attribute the overall
performance gain to its constituent sources. Table~\ref{tab:decompose}
summarizes the three-condition comparison.

\begin{table}[t]
\centering
\caption{Gain decomposition on 10~AIRS-Bench tasks (3~seeds, 50~steps,
invalid seeds scored as~0). Each row adds one capability layer.}
\label{tab:decompose}
\small
\begin{tabular}{@{}p{3.0cm}cccccr@{}}
\toprule
\textbf{Condition} & \textbf{Scaffold} & \textbf{Struct.\ state} &
\textbf{State evol.} & \textbf{Avg~NS} & \textbf{Valid~\%} &
\textbf{$\Delta$ vs prev.} \\
\midrule
AIRA-dojo   & \texttimes & \texttimes & \texttimes & 0.319 & 19\% & --- \\
ABL ($-\zt$) & \checkmark & \checkmark & \texttimes & 0.618 & 30\% & $+$0.299 \\
NAT ($+\zt$) & \checkmark & \checkmark & \checkmark & 0.616 & 40\% & $-$0.002 \\
\bottomrule
\end{tabular}
\end{table}

\paragraph{Scaffold gains dominate.}
The execution-aware scaffold (Section~\ref{sec:scaffold}) accounts for the
vast majority of the improvement: ABL raises Avg~NS from 0.319 to 0.618
($+$93.7\%). This gain is driven by two effects. First, environment-specific
prompt constraints eliminate entire categories of bugs---AIRA-dojo fails to
produce any valid submission on 3~of the 10~tasks (CvQm9,
SentimentAnalysis, CodeRetrieval), whereas ABL scores on all~10. Second,
environment reconnaissance gives every operator accurate knowledge of
available packages, GPU capabilities, and API compatibility, reducing
trial-and-error on environment-related issues.

\paragraph{Cognitive-state evolution: nuanced effect.}
NAT does not improve Avg~NS over ABL in aggregate ($-$0.002). However, NAT's
valid submission rate is 10~percentage points higher (40\% vs~30\%),
indicating that accumulated cognitive state helps the agent \emph{avoid
repeated errors} and produce more executable code. The lack of NS improvement
despite higher validity is explained by a \emph{direction-locking} effect
analyzed in Section~\ref{sec:boundary}: when the cognitive state accumulates
incorrect hypotheses early, the agent persists in a suboptimal direction
across the remaining steps, increasing cross-seed variance and offsetting
gains from better seeds.

\paragraph{Implication for evaluation.}
The clean mechanism comparison for cognitive-state evolution is therefore
ABL~vs~NAT, not Baseline~vs~NAT. The Baseline~$\to$~ABL gap reflects
engineering improvements that are shared infrastructure, not the cognitive
mechanism under study. All subsequent analyses use ABL as the primary
reference.


\subsection{Causal intervention analysis}
\label{sec:causal}

Using the four conditions defined in Table~\ref{tab:interventions}, we test
whether the performance gain from cognitive-state evolution is causal and
task-specific. All conditions use the Greedy topology.

\subsubsection{Aggregate results}

\begin{table}[t]
\centering
\caption{Causal intervention results (Greedy, 10~tasks $\times$ 3~seeds).
NAT: natural evolution. ABL: ablated (blank state each step).
SCR and FRZ conditions are pending completion.}
\label{tab:causal}
\begin{tabular}{@{}lccc@{}}
\toprule
\textbf{Condition} & \textbf{Avg~NS} & \textbf{Valid~\%} &
$\boldsymbol{\Delta}$\textbf{ vs NAT} \\
\midrule
NAT (natural)            & 0.616 & 40\% & ---     \\
ABL (ablated)            & 0.612 & 30\% & $-$0.005 \\
FRZ ($N_f{=}5$, frozen)  & \multicolumn{3}{c}{\emph{pending}} \\
SCR (scrambled)           & \multicolumn{3}{c}{\emph{pending}} \\
\bottomrule
\end{tabular}
\end{table}

The NAT--ABL comparison reveals a nuanced picture. While the Avg~NS difference
is small ($+$0.005), NAT's valid submission rate is 10~percentage points higher
(40\% vs 30\%), demonstrating that cognitive-state evolution helps the agent
\emph{avoid repeated errors} and produce more executable code. The best-score
gap is narrow because ABL's random exploration occasionally discovers strong
solutions that NAT's directed search misses---a trade-off between
\emph{consistency} (NAT) and \emph{diversity} (ABL).

\subsubsection{Per-task breakdown}

\begin{table}[t]
\centering
\caption{Per-task normalized scores: AIRA-dojo baseline, ABL, and NAT
(Greedy, 3-seed average, invalid seeds scored as~0).
Tasks ordered by NAT $-$ Baseline $\Delta$.}
\label{tab:pertask}
\small
\begin{tabular}{@{}lccccr@{}}
\toprule
\textbf{Task} & \textbf{Tier} & \textbf{Baseline} & \textbf{ABL} &
\textbf{NAT} & \textbf{NAT$-$ABL} \\
\midrule
SQuAD Exact Match       & Med  & 0.737 & 0.363 & 0.754 & $+$0.391 \\
Qm9 MAE                & Hard & 0.000 & 0.733 & 0.747 & $+$0.014 \\
Zinc MAE               & Hard & 0.196 & 0.681 & 0.698 & $+$0.017 \\
SVAMP Accuracy          & Exp  & 0.165 & 0.201 & 0.180 & $-$0.021 \\
SICK Accuracy           & Easy & 1.010 & 1.111 & 1.089 & $-$0.022 \\
Code Retrieval MRR      & Hard & 0.001 & 0.322 & 0.286 & $-$0.036 \\
Solar Weekly MAE        & Easy & 0.837 & 0.937 & 0.897 & $-$0.040 \\
Yelp Accuracy           & Med  & 0.000 & 0.744 & 0.677 & $-$0.068 \\
SuperGLUE WSC Accuracy  & Exp  & 0.206 & 0.313 & 0.219 & $-$0.094 \\
Winogrande Accuracy     & Med  & 0.040 & 0.711 & 0.617 & $-$0.094 \\
\midrule
\textbf{Average}        &      & 0.319 & 0.612 & 0.616 & $+$0.005 \\
\bottomrule
\end{tabular}
\end{table}

Table~\ref{tab:pertask} reveals strong task-dependent variation.
NAT outperforms ABL on SQuAD Exact Match by a large margin ($+$0.391),
consistent with the hypothesis that reading comprehension requires sustained
knowledge accumulation across steps. NAT also shows modest gains on
molecular property prediction tasks (Qm9, Zinc), where iterative refinement
of model architecture benefits from accumulated failure patterns.

Conversely, ABL outperforms NAT on several tasks, most notably Winogrande
($-$0.094) and WSC ($-$0.094). Analysis of per-seed variance
(Section~\ref{sec:boundary}) reveals that this is driven by NAT's
\emph{direction-locking} effect: when the cognitive state accumulates
incorrect hypotheses early on, the agent commits to a suboptimal approach
across the remaining steps. ABL, which resets to a blank state each step,
maintains higher exploration diversity and occasionally discovers better
solutions by chance.


\subsection{Analysis: when does cognitive state help?}
\label{sec:boundary}

We analyze the per-task intervention effects to identify the conditions under
which cognitive-state evolution provides the largest---and smallest---gains.

% [placeholder] The following three hypotheses will be evaluated against data.
% Rewrite each paragraph with specific numbers and task examples once
% Table~\ref{tab:pertask} is populated.

\paragraph{Finding~1: cognitive state helps most on knowledge-accumulation
tasks.}
SQuAD Exact Match shows the largest NAT--ABL gap ($+$0.391). This task
requires the agent to iteratively refine a reading comprehension pipeline
across multiple attempts, where insights about data formatting, model
selection, and training stability carry forward across steps. The cognitive
state enables NAT to build on prior diagnoses rather than rediscovering them.
Molecular property prediction tasks (Qm9 $+$0.014, Zinc $+$0.017) show
smaller but consistent positive effects for similar reasons.

\paragraph{Finding~2: cognitive state introduces direction-locking risk.}
On Winogrande ($-$0.094) and WSC ($-$0.094), NAT underperforms ABL.
Per-seed analysis reveals that NAT's standard deviation across seeds is
3--6$\times$ larger than ABL's (e.g., Winogrande NAT~std = 0.242 vs
ABL~std = 0.039). When the cognitive state locks onto an incorrect hypothesis
early, the agent persists in a suboptimal direction. ABL's stateless design
provides implicit regularization through exploration diversity.

\paragraph{Finding~3: ceiling effects limit gains on easy tasks.}
SICK Accuracy achieves NS~$>$~1.0 under both NAT (1.089) and ABL (1.111),
surpassing published SOTA. At this performance level, the cognitive state
provides no additional benefit---a single well-crafted fine-tuning pipeline
suffices. The small negative gap ($-$0.022) is within noise.

\paragraph{Finding~4: system-level scaffold is a prerequisite, not a
substitute.}
As shown in Section~\ref{sec:decompose}, the scaffold improvements account
for the bulk of the Baseline-to-NAT gain. However, the scaffold alone does
not produce cognitive accumulation---ABL's valid submission rate plateaus at
30\%, while NAT reaches 40\% through error avoidance learned via reflection.
The scaffold creates the conditions under which cognitive-state evolution
\emph{can} be effective; whether it \emph{is} effective depends on the state
management mechanism (Section~\ref{sec:inner}) and the task characteristics
discussed above.


\subsection{Cognitive state dynamics}
\label{sec:dynamics}

To understand \emph{how} the cognitive state evolves during a run, we analyze
the trajectory logs recorded by the instrumentation module.

\paragraph{IQ trajectory.}
% [placeholder] Plot IQ(z_t) vs step for representative tasks (one easy, one
% hard). Report: (1) rate of IQ growth, (2) whether IQ plateaus and at which
% step, (3) whether high-IQ trajectories correlate with better final NS.

\paragraph{State field evolution.}
% [placeholder] For Greedy NAT on a representative task, report:
% (1) number of hypotheses added/removed per step (from trajectory.jsonl),
% (2) learned patterns accumulation curve,
% (3) confidence trajectory vs actual success rate (calibration over time),
% (4) fraction of reflection calls that fell back to z_{t+1} = z_t.

\paragraph{Reflection cost.}
% [placeholder] Report:
% (1) average tokens consumed per reflect call,
% (2) reflect latency as fraction of total step time,
% (3) reflect fallback rate (JSON parse failure).


% ====================================================================
% Section 5: Discussion
% ====================================================================
\section{Discussion}
\label{sec:discussion}

\paragraph{System improvements as prerequisite.}
A large fraction of the Baseline-to-ABL gain (Section~\ref{sec:decompose})
comes from execution-aware system improvements rather than cognitive-state
evolution. This does not invalidate the cognitive-state contribution; rather,
it defines a stronger and more reliable substrate on top of which state
evolution is evaluated. The lesson is practical: before an agent can benefit
from \emph{thinking}, it must first be able to \emph{code}. Environment
reconnaissance, compatibility-aware prompting, and failure-preserving output
truncation address this prerequisite. The cleaner mechanism comparison is
therefore ABL~vs~NAT---both share the improved scaffold, and any difference
is attributable to cognitive-state accumulation alone.

\paragraph{From memory to cognition.}
Existing agent memory systems---skill libraries~\citep{wang2023voyager},
experience pools~\citep{zhao2024expel}, hierarchical
caches---store \emph{what happened}. DLE-Agent's cognitive state encodes
\emph{what the agent currently believes and should do next}. This distinction
mirrors the difference between episodic and working memory in cognitive
science. The causal intervention results (Section~\ref{sec:causal}) provide
initial evidence for this distinction; the planned scrambled-state (SCR)
condition will test whether generic text enrichment alone can explain the
observed effects.

\paragraph{Operator design vs.\ search topology.}
Consistent with AIRA-dojo~\citep{airadojo2025}, we observe that per-step
decision quality matters more than search-tree sophistication. The cognitive
state improves per-step quality by providing accumulated context, which
explains why even the Chain topology benefits from~$\zt$. MC-ESES further
amplifies this by searching over cognitive states rather than code solutions,
but the per-step improvement from reflection is the more fundamental
contribution.

% [placeholder] Support with specific numbers: Chain +z_t vs Greedy -z_t.

\paragraph{Cost of reflection.}
Each reflection step requires one additional LLM call.
% [placeholder] Report actual overhead from Section~\ref{sec:dynamics}:
% tokens per reflect, latency fraction, total API cost increase.
% Replace the speculative "15-20%" with measured values.

\paragraph{Limitations.}
Our evaluation uses a single backbone LLM (Qwen-3.5-397B). While the
cognitive-state mechanism is model-agnostic, absolute performance numbers may
differ with other models. The causal intervention experiments isolate the
cognitive state's contribution within a fixed search topology; cross-topology
interactions (e.g., whether MC-ESES branching amplifies or dampens the
reflection effect) warrant further study. The 50-step budget, while sufficient
to demonstrate the mechanism, may not capture long-horizon dynamics that emerge
over hundreds of steps. Finally, the full replacement update mechanism means
the LLM can inadvertently discard useful information during
reflection---append-only or hybrid update strategies are worth exploring.


% ====================================================================
% Section 6: Conclusion
% ====================================================================
\section{Conclusion}
\label{sec:conclusion}

We introduced DLE-Agent, a dual-level evolution framework that equips AI
research agents with an explicit, persistent cognitive state. By coupling
branch-level search (outer level) with structured cognitive-state refinement
(inner level), DLE-Agent transforms research from stateless trial-and-error
into a closed-loop process of exploration and learning. On AIRS-Bench,
cognitive-state evolution yields consistent improvements across three search
topologies. Causal intervention experiments confirm that the cognitive state
carries task-specific information that causally improves downstream decisions.

% [placeholder] Add 1-2 sentences with headline numbers once data is final.

Our work opens several directions: (i)~scaling to longer horizons and more
complex tasks; (ii)~learning the reflection operator from data rather than
prompting; (iii)~cross-task transfer of cognitive states as a form of
meta-learning; and (iv)~hybrid update mechanisms that combine the flexibility
of full replacement with the safety of append-only constraints.


\section*{Author Contributions}
% If you'd like to, you may include a section for author contributions as is done
% in many journals. This is optional and at the discretion of the authors.

\section*{Acknowledgments}
% Use unnumbered first level headings for the acknowledgments. All
% acknowledgments, including those to funding agencies, go at the end of the paper.

\section*{Ethics Statement}
This work develops autonomous AI research agents that generate and execute code
within sandboxed environments. All experiments run in isolated containers with
no network access, posing minimal security risk. We do not release trained
models or fine-tune on sensitive data. The benchmark tasks (AIRS-Bench) use
publicly available datasets. We acknowledge that increasingly capable AI
research agents raise broader questions about the future of scientific labor;
we believe transparent evaluation and mechanistic understanding---as pursued in
our causal intervention experiments---are essential steps toward responsible
deployment.


\bibliography{references}
\bibliographystyle{colm2026_conference}

\appendix

\section{Cognitive State Fields}
\label{sec:app_state}

Table~\ref{tab:state_fields} provides the complete specification of the
cognitive state $\zt$.

\begin{table}[ht]
\centering
\caption{Fields of the cognitive state $\zt$.}
\label{tab:state_fields}
\small
\begin{tabular}{@{}llp{5.5cm}@{}}
\toprule
\textbf{Field} & \textbf{Type} & \textbf{Description} \\
\midrule
$\tau$ (task\_understanding)  & string  & Structured interpretation of the
task, refined over time \\
$H$ (hypotheses)              & list    & Active working hypotheses about
promising approaches \\
$P$ (learned\_patterns)       & list    & Failure/success patterns extracted
from execution feedback \\
$S$ (attempt\_summaries)      & list    & Compressed summaries of prior
attempts (code sketch, result, lesson) \\
$D^{+}$ (preferred\_dirs)     & list    & Directions considered promising \\
$D^{-}$ (avoided\_dirs)       & list    & Directions ruled out as dead ends \\
$c$ (confidence)              & float   & Self-assessed confidence in current
best approach, $[0, 1]$ \\
$\varepsilon$ (env\_context)  & string  & Environment reconnaissance
(packages, GPU, known pitfalls) \\
\bottomrule
\end{tabular}
\end{table}


\section{Feedback Signal Dimensions}
\label{sec:app_feedback}

Table~\ref{tab:feedback_fields} describes the eight dimensions of the feedback
signal $\rt$.

\begin{table}[ht]
\centering
\caption{Dimensions of the multi-dimensional feedback signal $\rt$.}
\label{tab:feedback_fields}
\small
\begin{tabular}{@{}llp{5.5cm}@{}}
\toprule
\textbf{Dimension} & \textbf{Type} & \textbf{Description} \\
\midrule
metric          & float  & Raw evaluation score (None if buggy) \\
is\_buggy       & bool   & Whether code produced valid output \\
error\_category & string & Classified error type: environment,
logic, data, resource, none \\
error\_pattern  & string & Specific error signature for pattern
learning \\
trend           & string & Metric trajectory: improving,
stagnating, degrading, first \\
$\Delta$        & float  & Signed change vs.\ previous best \\
novelty         & float  & Jaccard distance from recent attempts,
$[0, 1]$ \\
$\Delta F$      & float  & Estimated free-energy reduction (proxy),
$[0, 1]$ \\
\bottomrule
\end{tabular}
\end{table}


\section{Operator Prompts}
\label{sec:app_prompts}

The cognitive state is rendered as a structured Markdown section and prepended
to all operator prompts (draft, improve, debug). The rendering includes:
\texttt{\#\# Task Understanding}, \texttt{\#\# Active Hypotheses},
\texttt{\#\# Learned Patterns}, \texttt{\#\# Attempt Summaries},
\texttt{\#\# Preferred Directions}, \texttt{\#\# Avoided Directions},
and \texttt{\#\# Confidence}. When
cross-branch knowledge transfer is active (MC-ESES), an additional
\texttt{\#\# Cross-Branch Insights} section is appended.

The reflect operator receives the current state $\zt$, the executed code, the
terminal output, the analysis summary, and the structured feedback $\rt$. It
produces a JSON object with updated values for all eight semantic fields plus a
\texttt{key\_insight} summary.


\section{Task Descriptions}
\label{sec:app_tasks}

Table~\ref{tab:tasks} lists the 10~tasks used in our experiments with their
domains, metrics, and difficulty tiers.

\begin{table}[ht]
\centering
\caption{Task selection for Exp~A and Exp~B.}
\label{tab:tasks}
\small
\begin{tabular}{@{}llll@{}}
\toprule
\textbf{Task} & \textbf{Domain} & \textbf{Metric} & \textbf{Tier} \\
\midrule
SICK Accuracy           & NLP         & Accuracy $\uparrow$   & Easy \\
Solar Weekly            & Time Series & MAE $\downarrow$      & Easy \\
Yelp Review Full        & NLP         & Accuracy $\uparrow$   & Medium \\
Winogrande              & NLP         & Accuracy $\uparrow$   & Medium \\
SQuAD                   & NLP         & Exact Match $\uparrow$& Medium \\
Qm9 (Cv)               & Chemistry   & MAE $\downarrow$      & Hard \\
ZINC                    & Chemistry   & MAE $\downarrow$      & Hard \\
CodeXGlue Retrieval     & Code        & MRR $\uparrow$        & Hard \\
SuperGLUE WSC           & NLP         & Accuracy $\uparrow$   & Expert \\
SVAMP                   & Math        & Accuracy $\uparrow$   & Expert \\
\bottomrule
\end{tabular}
\end{table}


\end{document}
